\documentclass[10pt,aspectratio=169]{beamer}

\usetheme{Madrid}
\usecolortheme{default}

\usepackage[T1]{fontenc}
\usepackage[utf8]{inputenc}
\usepackage[french]{babel}
\usepackage{booktabs}
\usepackage{array}
\usepackage{hyperref}

\title{CLI Options in the Wild}
\subtitle{Compter l'évolution des options en ligne de commande}
\author{
  Théophile BERTHELOT \and
  Loïs CHOW-CHINE-GALINAUD \and
  Yoan FRANÇOIS \and
  Baptiste GUERIF
}
\institute{Projet PROJ104 -- Encadrante : Xhevahire TËRNAVA}
\date{2025/2026}

\setbeamertemplate{navigation symbols}{}
\setbeamertemplate{footline}[frame number]

\begin{document}

\begin{frame}
  \titlepage
\end{frame}

\begin{frame}{Objectif du projet}
  Comprendre comment les interfaces en ligne de commande évoluent dans le temps.

  \vspace{0.4cm}
  L'idée centrale :
  \begin{itemize}
    \item récupérer l'aide de programmes CLI ;
    \item détecter et compter leurs options ;
    \item comparer les résultats entre programmes et entre versions ;
    \item visualiser l'évolution de cette complexité.
  \end{itemize}

  \vspace{0.2cm}
  Exemples : \texttt{ls}, \texttt{cat}, \texttt{grep}, \texttt{sed}, \texttt{sort}, \texttt{curl}, \texttt{git}, \texttt{coreutils}.
\end{frame}

\begin{frame}{Pourquoi les options CLI ?}
  Les options sont une mesure concrète de la richesse fonctionnelle d'un outil.

  \begin{itemize}
    \item Elles montrent l'ajout progressif de fonctionnalités.
    \item Elles révèlent la complexification des interfaces.
    \item Elles permettent de comparer outils Unix historiques et outils modernes.
    \item Elles exposent les limites de la documentation analysable automatiquement.
  \end{itemize}

  \vspace{0.4cm}
  \begin{block}{Question de fond}
    Est-ce que les outils en ligne de commande deviennent plus complexes au fil du temps ?
  \end{block}
\end{frame}

\begin{frame}{Fonctionnement général}
  Pipeline du projet :

  \begin{enumerate}
    \item choisir un programme ou un dépôt Git ;
    \item exécuter la commande d'aide adaptée ;
    \item extraire les options avec un parseur Python ;
    \item compter les options uniques ;
    \item exporter les résultats en CSV ou JSON ;
    \item produire des graphiques et visualisations.
  \end{enumerate}

  \vspace{0.3cm}
  Le projet contient aussi une interface graphique pour tester rapidement le compteur.
\end{frame}

\begin{frame}{Ce que le compteur analyse}
  Le parseur repère principalement :

  \begin{itemize}
    \item les options courtes : \texttt{-h}, \texttt{-v}, \texttt{-o} ;
    \item les options longues : \texttt{--help}, \texttt{--version}, \texttt{--output} ;
    \item les variantes d'aide selon les programmes.
  \end{itemize}

  \vspace{0.3cm}
  Modes de sortie :
  \begin{itemize}
    \item affichage terminal ;
    \item JSON ;
    \item CSV ;
    \item intégration possible avec des graphiques.
  \end{itemize}

  \vspace{0.2cm}
  \alert{Limite :} le projet extrait les options visibles dans l'aide, sans analyser toute la grammaire CLI.
\end{frame}

\begin{frame}{Difficulté : les aides ne sont pas uniformes}
  Tous les programmes ne documentent pas leurs options de la même manière.

  \begin{itemize}
    \item \texttt{--help} suffit pour certains programmes.
    \item D'autres nécessitent \texttt{-h}, \texttt{-help}, \texttt{--usage}, \texttt{--help all}, \texttt{--help full}.
    \item Certaines options sont groupées.
    \item Certaines aides contiennent des sous-commandes.
    \item Les formats changent parfois entre versions.
  \end{itemize}

  \vspace{0.3cm}
  Cela oblige à adapter la commande d'aide et les règles de parsing.
\end{frame}

\begin{frame}{Difficulté : voyager dans le temps}
  Pour mesurer l'évolution d'un programme, il faut reconstruire d'anciennes versions.

  \begin{enumerate}
    \item cloner un dépôt ;
    \item sélectionner les tags pertinents ;
    \item choisir une version par année ;
    \item compiler chaque version ;
    \item exécuter le compteur ;
    \item stocker les résultats.
  \end{enumerate}

  \vspace{0.3cm}
  \begin{alertblock}{Problème principal}
    Les commandes de build et les dépendances changent selon les versions.
  \end{alertblock}
\end{frame}

\begin{frame}{Enjeux techniques}
  Les principaux enjeux techniques sont :

  \begin{itemize}
    \item parsing robuste de textes d'aide hétérogènes ;
    \item prise en charge des alias et variantes d'aide ;
    \item gestion des sous-commandes ;
    \item automatisation Git sur plusieurs tags ;
    \item compilation de versions anciennes ;
    \item export propre des données ;
    \item visualisation lisible des résultats.
  \end{itemize}

  \vspace{0.3cm}
  Le projet combine parsing, automatisation système et analyse de données.
\end{frame}

\begin{frame}{Répartition des tâches}
  \small
  \begin{tabular}{p{0.34\linewidth} p{0.58\linewidth}}
    \toprule
    \textbf{Membre} & \textbf{Rôle principal} \\
    \midrule
    Yoan FRANÇOIS & Interface graphique et intégration des graphes \\
    Baptiste GUERIF & Code du compteur, parsing, restructuration du dépôt \\
    Loïs CHOW-CHINE-GALINAUD & Analyse des dépôts Git et automatisation temporelle \\
    Théophile BERTHELOT & Analyse de données, sources, graphiques temporels \\
    \bottomrule
  \end{tabular}

  \vspace{0.5cm}
  Travail commun :
  \begin{itemize}
    \item choix des règles de comptage ;
    \item tests sur programmes réels ;
    \item documentation ;
    \item préparation des livrables.
  \end{itemize}
\end{frame}

\begin{frame}{Avancement actuel}
  Ce qui est déjà présent dans le dépôt :

  \begin{itemize}
    \item compteur Python fonctionnel ;
    \item options de sortie JSON et CSV ;
    \item documentation utilisateur ;
    \item règles de comptage ;
    \item interface graphique Tkinter ;
    \item premiers graphiques ;
    \item script de téléchargement et test de versions Git ;
    \item planning et suivi de projet.
  \end{itemize}

  \vspace{0.3cm}
  Structure principale : \texttt{src/}, \texttt{docs/}, \texttt{Visualisation/}, \texttt{presentations/}.
\end{frame}

\begin{frame}{Planning de fin de projet}
  \begin{block}{Fin mai -- début juin}
    Corriger les cas d'erreur, documenter les limites, stabiliser le pipeline Git, sélectionner les programmes à analyser.
  \end{block}

  \begin{block}{Mi-juin}
    Lancer les analyses finales, générer les CSV/JSON, produire les graphiques Matplotlib, comparer l'évolution du nombre d'options.
  \end{block}

  \begin{block}{Dernière semaine}
    Finaliser README et documentation, préparer poster et présentation, répéter la soutenance.
  \end{block}
\end{frame}

\begin{frame}{Ce qu'il reste à faire}
  \begin{columns}[T]
    \begin{column}{0.5\linewidth}
      \textbf{Priorités techniques}
      \begin{itemize}
        \item améliorer le support des programmes complexes ;
        \item gérer \texttt{--help full} et \texttt{--help all} ;
        \item fiabiliser les builds historiques ;
        \item ajouter des recettes par année ou tag ;
        \item sauvegarder les résultats.
      \end{itemize}
    \end{column}
    \begin{column}{0.5\linewidth}
      \textbf{Priorités analyse}
      \begin{itemize}
        \item choisir les programmes finaux ;
        \item vérifier manuellement un échantillon ;
        \item produire les graphes finaux ;
        \item interpréter les tendances observées.
      \end{itemize}
    \end{column}
  \end{columns}
\end{frame}

\begin{frame}{Enjeux sociétaux}
  Les outils CLI restent centraux pour :

  \begin{itemize}
    \item l'administration système ;
    \item le développement logiciel ;
    \item la cybersécurité ;
    \item l'automatisation industrielle ;
    \item la recherche reproductible.
  \end{itemize}

  \vspace{0.3cm}
  Comprendre leur complexité aide à évaluer l'accessibilité, la charge cognitive, la maintenabilité et la qualité de la documentation technique.
\end{frame}

\begin{frame}{Enjeux environnementaux}
  Points négatifs possibles :

  \begin{itemize}
    \item compiler de nombreuses versions consomme du temps CPU ;
    \item télécharger des dépôts et dépendances multiplie les transferts réseau ;
    \item reconstruire des environnements anciens peut être coûteux.
  \end{itemize}

  \vspace{0.3cm}
  Pistes de réduction :
  \begin{itemize}
    \item limiter le nombre de programmes analysés ;
    \item choisir une version par année plutôt que tous les tags ;
    \item mettre en cache les dépôts et résultats ;
    \item éviter de relancer les builds inutiles ;
    \item documenter les échecs non reproductibles.
  \end{itemize}
\end{frame}

\begin{frame}{Apport du projet}
  Le projet produit :

  \begin{itemize}
    \item un outil de comptage réutilisable ;
    \item une méthode pour comparer des outils CLI ;
    \item des données exportables ;
    \item des visualisations ;
    \item une réflexion sur l'évolution de la complexité logicielle.
  \end{itemize}

  \vspace{0.3cm}
  Il relie histoire des logiciels, ergonomie des interfaces, automatisation et analyse de données.
\end{frame}

\begin{frame}{Conclusion}
  CLI Options in the Wild cherche à mesurer une évolution souvent invisible :

  \vspace{0.3cm}
  \begin{block}{}
    La croissance de la complexité des interfaces en ligne de commande.
  \end{block}

  \vspace{0.3cm}
  Le projet avance sur trois axes :
  \begin{itemize}
    \item compter correctement les options ;
    \item automatiser l'analyse dans le temps ;
    \item rendre les résultats compréhensibles par des graphes et une présentation claire.
  \end{itemize}

  \vspace{0.2cm}
  Objectif final : une analyse fiable, documentée et présentable de plusieurs programmes CLI réels.
\end{frame}

\end{document}
